% Part: first-order-logic
% Chapter: sequent-calculus
% Section: rules-and-proofs

\documentclass[../../../include/open-logic-section]{subfiles}

\begin{document}

\iftag{FOL}
      {\olfileid{fol}{seq}{rul}}
      {\olfileid{pl}{seq}{rul}}

\olsection{القواعد وال\usetoken{P}{derivation}}

فيما يأتي، لتدل~$\Gamma, \Delta, \Pi, \Lambda$ على متتاليات منتهية
من ال!!{sentence}s.

\begin{defn}[التتابعية]
\emph{التتابعية} تعبير من الصورة
\[
\Gamma \Sequent \Delta
\]
حيث~$\Gamma$ و$\Delta$ متتاليتان منتهيتان (وربما خاليتان) من
!!{sentence}s اللغة~$\Lang L$. وتسمى~$\Gamma$ \emph{المقدَّم}، بينما
تسمى~$\Delta$ \emph{التالي}.
\end{defn}

\begin{explain}
الفكرة الحدسية الكامنة وراء التتابعية هي: إذا تحققت جميع
ال!!{sentence}s في المقدَّم، تحققت واحدة على الأقل من ال!!{sentence}s في
التالي. أي إنه إذا كان $\Gamma = \tuple{!A_1, \dots, !A_m}$ و
$\Delta = \tuple{!B_1, \dots, !B_n}$، فإن $\Gamma \Sequent \Delta$
تتحقق إذا وفقط إذا تحققت
\[
(!A_1 \land \cdots \land !A_m) \lif (!B_1 \lor \cdots \lor
!B_n)
\]
وثمة حالتان خاصتان: أن تكون~$\Gamma$ خالية، وأن تكون~$\Delta$ خالية.
فإذا كانت~$\Gamma$ خالية، أي إذا كان $m = 0$، فإن $\quad
\Sequent \Delta$ تتحقق إذا وفقط إذا تحقق
$!B_1 \lor \dots \lor !B_n$. وإذا كانت~$\Delta$ خالية، أي إذا كان
$n = 0$، فإن $\Gamma \Sequent \quad$ تتحقق إذا وفقط إذا تحقق
$\lnot(!A_1 \land \dots \land !A_m)$. ونقول إن التتابعية صالحة إذا
وفقط إذا كانت ال!!{sentence} المقابلة صالحة.
\end{explain}

إذا كانت~$\Gamma$ متتالية من ال!!{sentence}s، كتبنا $\Gamma, !A$
لنتيجة إلحاق~$!A$ بالطرف الأيمن من~$\Gamma$ (وكتبنا $!A, \Gamma$
لنتيجة إلحاق~$!A$ بالطرف الأيسر من~$\Gamma$). وإذا كانت~$\Delta$
متتالية من ال!!{sentence}s أيضًا، كانت~$\Gamma, \Delta$ وصل
المتتاليتين.

\begin{defn}[تتابعية ابتدائية]
\emph{التتابعية الابتدائية} تتابعية
\iftag{prvFalse,prvTrue}{من إحدى الصور الآتية:
  \begin{enumerate}
    \item $!A \Sequent !A$
    \tagitem{prvTrue}{$\quad \Sequent \ltrue$}{}
    \tagitem{prvFalse}{$\lfalse \Sequent \quad$}{}
  \end{enumerate}}
{من الصورة $!A \Sequent !A$} لأي !!{sentence}~$!A$ في اللغة.
\end{defn}

ال!!^{derivation}s في حساب التتابعيات أشجار معينة من التتابعيات؛ تكون
التتابعيات العليا فيها تتابعيات ابتدائية، وإذا وقعت تتابعية تحت
تتابعية أخرى أو تتابعيتين، وجب أن تتبع منها أو منهما على الوجه الصحيح
بقاعدة استدلال. وتنقسم قواعد~$\Log{LK}$ إلى نوعين رئيسين: قواعد
\emph{منطقية} وقواعد \emph{بنيوية}. وتسمى القواعد المنطقية باسم
!!{main operator} لل!!{sentence} التي تحتوي على~$!A$ و/أو~$!B$ في
التتابعية السفلى. ولكل قاعدة نسختان: إحداهما لاستنتاج تتابعية تقع فيها
ال!!{sentence} المحتوية على ال!!{operator} في الجانب الأيسر، والأخرى
لاستنتاج تتابعية تقع فيها تلك ال!!{sentence} في الجانب الأيمن.

\end{document}
